\chapter{Literature Review}
\label{ch2-LiteratureReview}

Using blockchain to trace goods through a supply chain is not a new idea.
Over the last few years researchers have tried it in a number of areas,
from food and medicine to forestry and factory manufacturing. We went
through ten of these earlier works before we started building
\emph{TimberTrace}. Some are published research papers and some are
open-source projects that other developers have put on GitHub. We did not
just read them for background; in each one we looked for a particular idea
worth borrowing. Sometimes we used the idea more or less as it was, other
times we changed it to fit timber, and in a few cases we decided a simpler
option made more sense for a student prototype. The rest of this chapter
goes through these works one by one and points out exactly which design
choice in our own project each of them shaped.

\section{Survey of Related Work}

The work closest to ours is by Figorilli et al.~\cite{figorilli2018wood}.
As far as we know, it was the first published attempt to trace wood
electronically across the whole supply chain using a blockchain. They put
RFID tags on the logs and wrote each scan to an immutable ledger, so the
physical wood and its digital record stayed linked. Westerkamp et
al.~\cite{westerkamp2018token} looked at the problem from a different
angle. In their ``token-recipe'' model, making a product means using up
the tokens that stand for its raw materials and creating a new token for
the finished item, which keeps the link between inputs and outputs on the
chain itself. Two more projects helped shape the consumer-facing side of
our work: the SCTS system~\cite{scts} let anyone check a product's history
by scanning a QR code, and the \verb|trace| dApp~\cite{trace} went a step
further by attaching signed certificates from trusted experts to each
product. A third, \verb|smart-traceability|~\cite{smarttraceability},
showed a neat way to handle large files by keeping the batch data on IPFS
and storing only its hash on-chain.

A second group of works deals mainly with who is allowed to do what.
PharmaChain~\cite{pharmachain} and the dual-contract design in
Ref.~\cite{dualcontract} both handle roles carefully; the latter even
splits the data and the permission logic into separate contracts so the
rules can be updated without disturbing the stored records. The dApp in
Ref.~\cite{faizack} treats the supply chain as a set of fixed stages that
a product can only move forward through, with role checks guarding each
step. Finally, two studies made us think harder about which blockchain to
build on. The Hyperledger Fabric system in Ref.~\cite{hyperledger} showed
what a permissioned, enterprise-grade network can do, and the recent
timber-focused framework in Ref.~\cite{timberSC2025} added something we
had overlooked, treating the transfer of ownership as its own recorded
action on the chain rather than just another event.

\section{Comparative Summary}

Table~\ref{tab:litreview} consolidates the survey, listing each
paper/project alongside the specific idea we carried into
\emph{TimberTrace}.

\renewcommand{\arraystretch}{1.35}
\begin{longtable}{|c|p{5.1cm}|p{6.6cm}|}
  \caption{Literature reviewed and the idea adopted in this project.}
  \label{tab:litreview} \\
  \hline
  \textbf{Sl.\ No.} & \textbf{Paper / Project Name} & \textbf{Idea Used in My Project} \\
  \hline
  \endfirsthead

  \multicolumn{3}{c}%
  {\tablename\ \thetable\ -- \textit{continued from previous page}} \\
  \hline
  \textbf{Sl.\ No.} & \textbf{Paper / Project Name} & \textbf{Idea Used in My Project} \\
  \hline
  \endhead

  \hline
  \multicolumn{3}{r}{\textit{continued on next page}} \\
  \endfoot

  \hline
  \endlastfoot

  1 & Figorilli et al. (2018), ``A Blockchain Implementation Prototype
  for the Electronic Open Source Traceability of Wood along the Whole
  Supply Chain,'' \emph{Sensors}, MDPI~\cite{figorilli2018wood} &
  Recording every physical handoff of the wood (harvest $\rightarrow$
  lot $\rightarrow$ batch $\rightarrow$ furniture) as an immutable,
  tamper-proof on-chain event, so the ledger is a permanent provenance
  record of each stage. \\
  \hline

  2 & Westerkamp, Victor \& K\"upper (2018), ``Blockchain-based Supply
  Chain Traceability: Token Recipes Model Manufacturing
  Processes,'' IEEE Blockchain~\cite{westerkamp2018token} &
  The input-consumed-to-output manufacturing model, realised as our
  UTXO-inspired mass-balance rule in which a Lot is consumed to create a
  Batch and a Batch to create Furniture, preserving the input--output
  link on-chain. \\
  \hline

  3 & SCTS --- Supply Chain Tracking System (AtrauraBlockchain,
  GitHub)~\cite{scts} &
  Placing a QR code on the physical product that resolves its full
  on-chain history, enabling anyone to audit provenance without special
  access --- our \texttt{/verify} QR flow. \\
  \hline

  4 & \texttt{maximevaillancourt/trace} --- certification dApp
  (GitHub)~\cite{trace} &
  Attaching signed certificates to products; we store a certificate
  IPFS CID inside the Lot and Furniture records as a lightweight
  certification layer over the provenance data. \\
  \hline

  5 & \texttt{ecobioca/smart-traceability} (GitHub)~\cite{smarttraceability} &
  Storing all bulky batch/product metadata (certificates, product
  images) off-chain on IPFS via Pinata and keeping only the content
  hash on-chain to save storage and gas. \\
  \hline

  6 & PharmaChain (Bapatla et al., 2022), \emph{IET
  Networks}~\cite{pharmachain} &
  Role-based access control enforced through Solidity function
  modifiers --- our \texttt{AUCTION\_HOUSE\_ROLE} and
  \texttt{FACTORY\_ROLE} guarded by OpenZeppelin's
  \texttt{onlyRole()}, with the admin granting/revoking roles. \\
  \hline

  7 & Dual-Contract Architecture with RBAC (2025), \emph{Scientific
  Reports}~\cite{dualcontract} &
  Separation of concerns between authoritative logic and data ---
  reflected in our decoupling of the authoritative smart contract (source
  of truth) from the read-optimised PostgreSQL mirror, and noted as a
  future dual-contract upgrade path. \\
  \hline

  8 & \texttt{faizack/Supply-Chain-Blockchain} (GitHub)~\cite{faizack} &
  Modelling the supply chain as defined, role-gated stages with a
  dedicated contract function for each transition
  (\texttt{registerLot}, \texttt{purchaseWood},
  \texttt{createFurniture}). \\
  \hline

  9 & High-Efficiency Blockchain-Based Supply Chain Traceability
  (Hyperledger Fabric, 2022)~\cite{hyperledger} &
  Throughput comparison of permissioned vs.\ public ledgers, which
  informed our deliberate choice of public Ethereum Sepolia for the
  prototype and identified Hyperledger as the future scaling path. \\
  \hline

  10 & Blockchain-Enabled Smart Contracts for Secure Timber
  Traceability (2025), ScienceDirect~\cite{timberSC2025} &
  On-chain ownership tracking (the \texttt{owner} field on each Lot,
  Batch and Furniture) together with QR-based consumer verification of
  the harvested-and-certified timber. \\
  \hline

\end{longtable}
\renewcommand{\arraystretch}{1.0}

\section{Research Gap}

Across the surveyed works, three consistent trade-offs emerge. First,
hardware-driven systems~\cite{figorilli2018wood} achieve automation at the
cost of expensive IoT infrastructure, whereas a software-only,
operator-signed model is more appropriate for a research prototype.
Second, token-heavy representations~\cite{westerkamp2018token,
smarttraceability} and per-product contracts~\cite{scts} offer flexibility
but at higher gas cost and reduced auditability compared with a single
monolithic contract using struct linkage. Third, enterprise permissioned
ledgers~\cite{hyperledger} provide superior throughput but demand
self-operated infrastructure. \emph{TimberTrace} occupies the gap between
these extremes: a software-only, single-contract Ethereum system that
enforces a mass-balance conservation invariant on-chain, mirrors state to
a fast read database off-chain, and exposes QR-based public verification
--- combining the most relevant ideas from the literature into a
lightweight, auditable timber-provenance prototype.


%%% Local Variables:
%%% mode: latex
%%% TeX-master: "../mainrep"
%%% End:
